\documentclass[12pt]{article}

\usepackage[margin=1in]{geometry}
\usepackage{amsmath, amssymb}
\usepackage{booktabs}
\usepackage{array}
\usepackage{hyperref}
\usepackage{natbib}
\usepackage{graphicx}
\usepackage{setspace}
\usepackage{xcolor}
\usepackage{caption}
\usepackage{lscape}
\usepackage{longtable}

\onehalfspacing

\title{\textbf{Methods: Estimating Infectious Disease Importation Risk\\
during the 2026 FIFA World Cup}}
\author{Jose Herrera-Diestra}
\date{\today}

\begin{document}

\maketitle

% ============================================================
\section*{Methods}
% ============================================================

\subsection*{Study setting}

The 2026 FIFA World Cup will be hosted jointly by the United States,
Mexico, and Canada, with matches scheduled from June~11 through
July~19, 2026. The United States will host the majority of matches
across eleven venues: MetLife Stadium (East Rutherford, NJ), SoFi
Stadium (Inglewood, CA), AT\&T Stadium (Arlington, TX), Mercedes-Benz
Stadium (Atlanta, GA), NRG Stadium (Houston, TX), Arrowhead Stadium
(Kansas City, MO), Lincoln Financial Field (Philadelphia, PA), Hard
Rock Stadium (Miami Gardens, FL), Levi's Stadium (Santa Clara, CA),
Lumen Field (Seattle, WA), and Gillette Stadium (Foxborough, MA).
All 48 qualified nations are listed in Supplementary Table~1.

We estimated the importation risk of four infectious diseases---dengue
fever, malaria, measles, and pertussis---for each host city using three
complementary models of increasing resolution: (1) a \emph{baseline
model} using historical regional air travel patterns (2023--2025);
(2) a \emph{World Cup (WC)-adjusted model} that upgrades travel volumes
to country-level 2026 projections; and (3) a \emph{schedule-driven
model} that separates World Cup fan travel from background tourism and
routes fans to the specific venue cities where their national team plays.

% ============================================================
\subsection*{Importation risk framework}

We modelled disease importation as a Poisson process. If
$\lambda_{h,d}$ is the expected number of imported cases of disease
$d$ arriving in city $h$ during June 2026, the probability of
observing at least one importation is:

\begin{equation}
  P(X_{h,d} \geq 1) = 1 - e^{-\Lambda_{h,d}}
  \label{eq:poisson}
\end{equation}

where $\Lambda_{h,d} = \sum_{s} \lambda_{s,h,d}$ sums over all origin
units $s$ (regions in the baseline model; countries in the
WC-adjusted model). This framework is consistent with previously
published importation risk analyses for large international sporting
events~\citep{Herrera2024}.

% ============================================================
\subsection*{Baseline model: regional travel volume (2023--2025)}

\subsubsection*{Travel data}

Origin-to-city air arrival volumes were obtained from the
US International Air Travel Statistics (I-92 programme, Bureau of
Transportation Statistics / US Department of
Commerce).\footnote{\url{https://www.trade.gov/us-international-air-travel-statistics-i-92-data}}
These data record monthly passenger movements between broad
world-origin regions and specific US gateway cities (Boston, Dallas,
Houston, Newark, New York, Philadelphia). Origin regions were used as
provided by the I-92 programme, with two modifications applied to
align with the country-region correspondence table derived from I-94
data: (1) Western Europe and Eastern Europe were merged into a single
Europe region; and (2) Mexico and Canada were retained as standalone
regions rather than being grouped into North America, given their
status as co-host nations and the scale of their visitor volumes.
Oceania was excluded from the routing fraction analysis due to the
absence of direct Oceania-to-US gateway city data in the I-92 files
used; Australian visitors are captured in the WC-adjusted model via
the NTTO Tier~1 projection.

Total monthly arrivals were defined as:
\begin{equation}
  N_{\text{arrivals}} = N_{\text{foreign originating}}
                      + N_{\text{foreign returning}}
                      + N_{\text{US citizen returning}}
\end{equation}

The mean June arrivals over 2023--2025 were used as the baseline
travel volume $N_{r,h}$ for region $r$ to city $h$, chosen to
capture the most recent pre-tournament summer travel patterns.

\subsubsection*{City routing fractions}

For the WC-adjusted model (described below), city routing fractions
were estimated from the I-92 data as:
\begin{equation}
  f_{r,h} = \frac{N_{r,h}}{N_{r,\text{US total}}}
\end{equation}
where $N_{r,\text{US total}}$ is the mean June arrivals from region
$r$ to the entire United States. These fractions represent the
share of arrivals from a given origin region that are routed to each
specific US gateway city.

\subsubsection*{Regional disease incidence}

For each disease $d$ and origin region $r$, a regional incidence
proxy was computed by aggregating country-level disease data within
the region:

\begin{equation}
  I_{r,d} = \rho_d \cdot \frac{C_{r,d}}{P_r}
  \label{eq:regional_inc}
\end{equation}

where $C_{r,d}$ is the total reported cases for disease $d$ in
region $r$ during June (or the annual equivalent divided by 12 for
annual data), $P_r$ is the total population of region $r$, and
$\rho_d$ is a disease-specific adjustment factor that accounts for
reporting completeness and the fraction of reported cases that are
actively infectious (see Table~\ref{tab:params}).

\subsubsection*{Importation intensity (baseline)}

The expected importation intensity from region $r$ to city $h$ for
disease $d$ is:

\begin{equation}
  \lambda_{r,h,d} = N_{r,h} \cdot I_{r,d} \cdot p_d
  \label{eq:lambda_baseline}
\end{equation}

where $p_d$ is the probability that an infected individual from
region $r$ travels internationally while infectious. The total
importation intensity for city $h$ is then:

\begin{equation}
  \Lambda_{h,d} = \sum_{r} \lambda_{r,h,d}
\end{equation}

% ============================================================
\subsection*{World Cup--adjusted model: country-level travel
volume (2026 projections)}

\subsubsection*{Motivation}

The baseline model uses historical regional air travel data that
reflects normal summer travel patterns and does not account for the
additional influx of visitors attending the World Cup. Furthermore, I-92
data aggregate countries into broad regions, masking the
heterogeneity in disease burden between countries within the same
region. The WC-adjusted model addresses both limitations.

\subsubsection*{Construction of 2026 travel volume}

Country-level June 2026 visitor volumes to the United States were
constructed from two complementary sources using a tiered approach:

\paragraph{Tier 1 --- NTTO projections (top 12 source markets).}
Official 2026 international visitor projections for twelve major
source markets were obtained from the US National Travel and Tourism
Office (NTTO) 2025 Forecast
Report.\footnote{\url{https://www.trade.gov/travel-and-tourism-forecasts}}
These projections---available for Australia, Brazil, Canada, China,
France, Germany, India, Italy, Japan, Mexico, South Korea, and the
United Kingdom---already incorporate the World Cup tourism effect and
should \emph{not} have additional WC visitors added on top. The
2024-to-2026 growth factor for each country was derived as
$\phi_c = V_{c,2026} / V_{c,2024}$, where $V_{c,y}$ is the
projected annual visitor volume in year $y$ (Table~\ref{tab:ntto}).

\paragraph{Tier 2 --- COR data scaled by global growth factors (remaining countries).}
For all other countries, country-level June 2024 arrivals were
extracted from the US Customs and Border Protection (CBP) I-94
Monthly Arrivals by Country of Residence
dataset.\footnote{\url{https://travel.trade.gov}} A growth factor was
then assigned based on World Cup qualification status:

\begin{itemize}
  \item \emph{WC-qualified countries not in Tier 1}: scaled by the
    overall NTTO 2024$\to$2026 growth factor ($\phi_{\text{WC}} =
    85{,}017/72{,}390 = 1.174$), which incorporates the WC tourism
    uplift at the aggregate level.
  \item \emph{Non-qualified countries}: scaled by the estimated
    baseline growth factor without the WC effect
    ($\phi_{\text{base}} = 1.134$, corresponding to approximately
    $6.5\%$ annual growth over two years).
\end{itemize}

The projected June 2026 arrivals for country $c$ are thus:
\begin{equation}
  N_{c,\text{June 2026}} = N_{c,\text{June 2024}}^{\text{COR}}
  \cdot \phi_c
  \label{eq:travel_wc}
\end{equation}

\paragraph{Country-to-city allocation.}
City-level arrival volumes for country $c$ and city $h$ were
estimated by applying the I-92 routing fractions from the baseline
model:
\begin{equation}
  N_{c,h}^{2026} = N_{c,\text{June 2026}} \cdot f_{r(c),h}
  \label{eq:routing}
\end{equation}
where $r(c)$ denotes the broad world region to which country $c$
belongs.

\subsubsection*{Country-level disease incidence}

In the WC-adjusted model, disease burden is computed at the country
level rather than the regional level, providing greater precision:

\begin{equation}
  I_{c,d} = \rho_d \cdot \frac{C_{c,d}}{P_c}
  \label{eq:country_inc}
\end{equation}

where $C_{c,d}$ is the reported cases of disease $d$ in country $c$
and $P_c$ is the national population. The same $\rho_d$ values as
in the baseline model are used to maintain comparability.

\subsubsection*{Importation intensity (WC-adjusted)}

\begin{equation}
  \lambda_{c,h,d} = N_{c,h}^{2026} \cdot I_{c,d} \cdot p_d
  \label{eq:lambda_wc}
\end{equation}

\begin{equation}
  \Lambda_{h,d} = \sum_{c} \lambda_{c,h,d}
\end{equation}

The probability of at least one importation is then computed via
Equation~\ref{eq:poisson}.

% ============================================================
\subsection*{Schedule-driven venue routing model}

\subsubsection*{Motivation}

Both the baseline and WC-adjusted models allocate international arrivals
to US gateway cities using I-92 routing fractions, which reflect
historical tourist flow patterns. However, World Cup fans travel to the
specific cities where their national team plays --- a fundamentally
different spatial logic. A supporter from Brazil attending a group-stage
match in Houston is unlikely to distribute their travel to Boston or
Philadelphia in the same proportions as a regular tourist. The
schedule-driven model captures this distinction by decomposing total
travel into a WC-specific fan stream and a background tourist stream,
routing each according to a different spatial rule.

\subsubsection*{Travel decomposition}

For each source country $c$, total projected June 2026 arrivals
$N_{c,\text{June 2026}}$ (Equation~\ref{eq:travel_wc}) are decomposed
into two components:

\begin{itemize}
  \item \textit{WC-fan stream}: the marginal arrivals attributable to the
    World Cup, defined as the increment above the no-WC counterfactual:
    \begin{equation}
      N_c^{\text{WC}} = N_{c,\text{June 2024}}^{\text{COR}} \cdot
        \max\!\bigl(0,\; \phi_c - \phi_{\text{base}}\bigr)
      \label{eq:wc_fan_stream}
    \end{equation}

  \item \textit{Background stream}: arrivals that would have occurred
    regardless of the World Cup:
    \begin{equation}
      N_c^{\text{bg}} = N_{c,\text{June 2026}} - N_c^{\text{WC}}
      \label{eq:bg_stream}
    \end{equation}
\end{itemize}

For countries with NTTO growth factors below the baseline
($\phi_c < \phi_{\text{base}} = 1.134$, e.g., the United Kingdom at
$\phi = 1.107$), the WC increment is set to zero and all travel is
treated as background. For a WC-qualified country with the global
WC growth factor ($\phi_{\text{WC}} = 1.174$), the WC fan fraction
is approximately $3.4\%$ of projected arrivals:
$(1.174 - 1.134)/1.174 \approx 0.034$.

\subsubsection*{Schedule-based fan routing}

The 2026 FIFA World Cup match schedule was extracted from event records
for all 48 qualifying nations. Matches with at least one undetermined opponent
(labelled ``TBD'' in the schedule) are excluded from the fan routing
calculation; both group-stage and knockout-round entries with TBD
opponents are dropped. For each WC-qualified country $c$ with team $t_c$, a
schedule routing fraction is defined as:

\begin{equation}
  \omega_{c,v} = \frac{g_{c,v}}{G_c}
  \label{eq:schedule_routing}
\end{equation}

where $g_{c,v}$ is the number of matches played by team $t_c$ at venue
city $v$ and $G_c = \sum_v g_{c,v}$ is the total number of (known)
group-stage matches. This fraction represents the share of WC fans from
country $c$ expected to attend games in venue city $v$.

WC-fan arrivals at each venue city are then:
\begin{equation}
  N_{c,v}^{\text{WC}} = N_c^{\text{WC}} \cdot \omega_{c,v}
  \label{eq:wc_arrivals_venue}
\end{equation}

Background arrivals continue to use I-92 regional routing fractions
applied to all six gateway cities:
\begin{equation}
  N_{c,h}^{\text{bg}} = N_c^{\text{bg}} \cdot f_{r(c),h}
  \label{eq:bg_arrivals}
\end{equation}

where $r(c)$ is the world region of country $c$ and $f_{r,h}$ is the
routing fraction from Equation~\ref{eq:routing}.

\paragraph{Note on multi-team countries.}
England and Scotland are distinct World Cup participants but share a
single country-of-residence entry (``United Kingdom'') in the COR
arrivals data. Both teams' match venues are combined under this single
entry, and the routing fractions $\omega_{\text{UK},v}$ are computed
over the union of their games. Because the fractions sum to one, no
double-counting of UK arrivals occurs.

\paragraph{Venue cities.}
The schedule-driven model covers all sixteen host venues across the
three co-host nations (Table~\ref{tab:venues}). For cities that also
appear as I-92 gateway cities (Boston, Dallas, Houston, New York,
Philadelphia), the WC-fan and background streams are summed before
computing the importation intensity. For the remaining venue cities
(Atlanta, Kansas City, Los Angeles, Miami, San Francisco, Seattle in
the US; Toronto and Vancouver in Canada; Guadalajara, Mexico City, and
Monterrey in Mexico), only the WC-fan stream contributes --- background
tourist flows to those cities are not captured in the I-92 gateway data.

\begin{table}[h!]
\centering
\caption{WC 2026 host venues and their canonical city names used in the
  schedule-driven model. Cities also covered by I-92 gateway data are
  marked with $\dagger$.}
\label{tab:venues}
\begin{tabular}{llll}
\toprule
\textbf{Stadium} & \textbf{City (canonical)} & \textbf{Country} &
\textbf{I-92?} \\
\midrule
MetLife Stadium        & New York$^\dagger$    & USA    & Yes \\
SoFi Stadium           & Los Angeles           & USA    & No  \\
AT\&T Stadium          & Dallas$^\dagger$      & USA    & Yes \\
Mercedes-Benz Stadium  & Atlanta               & USA    & No  \\
NRG Stadium            & Houston$^\dagger$     & USA    & Yes \\
Arrowhead Stadium      & Kansas City           & USA    & No  \\
Lincoln Financial Field& Philadelphia$^\dagger$& USA    & Yes \\
Hard Rock Stadium      & Miami                 & USA    & No  \\
Levi's Stadium         & San Francisco         & USA    & No  \\
Lumen Field            & Seattle               & USA    & No  \\
Gillette Stadium       & Boston$^\dagger$      & USA    & Yes \\
BC Place               & Vancouver             & Canada & No  \\
BMO Field              & Toronto               & Canada & No  \\
Estadio Banorte        & Mexico City           & Mexico & No  \\
Estadio BBVA           & Monterrey             & Mexico & No  \\
Estadio Akron          & Guadalajara           & Mexico & No  \\
\bottomrule
\end{tabular}
\end{table}

\subsubsection*{Importation intensity (schedule-driven)}

The total Poisson importation intensity for venue city $v$ is the sum
of contributions from both travel streams across all source countries:

\begin{equation}
  \Lambda_{v,d} =
    \underbrace{%
      \sum_{c \in \mathcal{Q}}
        N_{c,v}^{\text{WC}} \cdot I_{c,d} \cdot p_d
    }_{\text{WC-fan stream}}
    \;+\;
    \underbrace{%
      \sum_{c}
        N_{c,h(v)}^{\text{bg}} \cdot I_{c,d} \cdot p_d
    }_{\text{background stream}}
  \label{eq:lambda_schedule}
\end{equation}

where $\mathcal{Q}$ denotes the set of WC-qualified countries with a
known group-stage schedule, and $h(v)$ maps venue city $v$ to its
corresponding I-92 gateway city (if one exists). The probability of at
least one importation is computed via Equation~\ref{eq:poisson}.

% ============================================================
\subsection*{Three-model comparison}

To quantify the effect of each modelling assumption, results from the
three models are compared side by side for each disease and destination
city. The three models form a nested hierarchy:

\begin{enumerate}
  \item \textbf{Baseline} (Sections above): I-92 regional travel
    (2023--2025 mean), regional disease incidence, five US I-92 gateway
    cities.
  \item \textbf{WC-adjusted}: Country-level 2026 travel projections
    (NTTO + COR scaling), country-level disease incidence, five US
    I-92 gateway cities (routing unchanged).
  \item \textbf{Schedule-driven}: Same 2026 projections decomposed into
    WC-fan and background streams; WC fans routed to all sixteen venue
    cities via the match schedule; background travelers use I-92 routing.
\end{enumerate}

The step from Model~1 to Model~2 isolates the combined effect of the
WC travel surge and country-level disease incidence resolution. The step
from Model~2 to Model~3 isolates the effect of schedule-based venue
routing: new cities (Atlanta, Miami, Kansas City, Los Angeles, San
Francisco, Seattle, and the three Mexican cities) appear in the
importation risk landscape only in Model~3.

Two summary metrics are compared (Figures~\ref{fig:comp_prob} and
\ref{fig:comp_lambda}):
\begin{itemize}
  \item $P(\geq 1)$: probability of at least one importation
    (Equation~\ref{eq:poisson}). Useful as a communication metric but
    saturates quickly toward~1 for high-risk cities.
  \item $\lambda$: expected number of imported cases (the raw Poisson
    intensity). Remains on a linear scale even when $P(\geq 1) \to 1$
    and allows direct quantitative comparison across cities and models.
\end{itemize}

% ============================================================
\subsection*{Country-level importation contributions}

The additive structure of the Poisson model allows city-level
importation risk to be attributed to individual source countries.
Total importation intensity decomposes as:

\begin{equation}
  \Lambda_{v,d} = \sum_c \lambda_{c,v,d}
  \label{eq:decompose}
\end{equation}

where $\lambda_{c,v,d} = N_{c,v} \cdot I_{c,d} \cdot p_d$ is the
contribution from source country $c$ to destination city $v$ for
disease $d$. Using the schedule-driven model (which separates WC-fan
and background streams), three summaries are produced:

\begin{enumerate}
  \item \textbf{Country importation ranking}: countries are ranked by
    total expected importations summed across all destination cities,
    $\sum_v \lambda_{c,v,d}$, separately for each disease
    (Figure~\ref{fig:country_rank}).

  \item \textbf{Travel-stream decomposition}: for the top 20
    countries by dengue importation burden, the fraction of
    $\lambda_{c,v,\text{dengue}}$ attributable to the WC-fan stream
    versus the background stream is shown
    (Figure~\ref{fig:stream_breakdown}). Countries with WC teams are
    expected to show a non-zero WC-fan component; non-qualified
    countries will show only background travel contributions.

  \item \textbf{Country $\times$ city importation heatmap}: a matrix
    visualisation of $\lambda_{c,v,d}$ summed over both streams,
    with rows ordered by total importation burden and columns by
    destination city (Figures~\ref{fig:heatmap_dengue}--\ref{fig:heatmap_pertussis}).
    This reveals the spatial concentration of risk: WC-qualified
    countries show importation concentrated in their specific match
    venues, while non-qualified countries show diffuse contributions
    spread across the five I-92 gateway cities.
\end{enumerate}

% ============================================================
\subsection*{Disease data sources}

\subsubsection*{Dengue}

Monthly country-level dengue case counts were obtained from the WHO
Global Dengue Surveillance dataset (accessed December 2025). For the
baseline model, cases were aggregated to broad world regions and the
mean June case count over 2024--2025 was used. For the WC-adjusted
model, the mean June 2024--2025 case count per country was used
directly, divided by national population to obtain a per-capita
incidence. Countries with no data were excluded.

\subsubsection*{Malaria}

Country-level malaria incidence rates for 2024 (cases per 1,000
population) were obtained from the WHO Global Malaria Programme
National Unit Data. Annual incidence was divided by 12 to obtain an
approximate monthly rate. Countries were recoded to match naming
conventions in the travel datasets where necessary (e.g.,
``Democratic Republic of the Congo'' $\to$ ``Zaire (formerly DRC)'').

\subsubsection*{Measles and pertussis}

Annual country-level incidence rates (cases per 1,000,000 population)
for measles and pertussis were obtained from the WHO Immunization Data
portal (accessed September and December 2025, respectively). Annual
incidence was divided by 12 to obtain an approximate monthly rate.

% ============================================================
\subsection*{Model parameters}

Table~\ref{tab:params} summarises the disease-specific parameters
used in all three models. The adjustment factor $\rho_d$ captures the
combined effect of under-reporting in routine surveillance and the
fraction of reported cases that represent actively infectious
individuals. The travel infectiousness probability $p_d$ reflects the
likelihood that a person with the disease continues to travel
internationally while infectious, accounting for the typical severity
and duration of symptoms relative to the incubation period.

\begin{table}[h!]
\centering
\caption{Disease-specific model parameters.}
\label{tab:params}
\begin{tabular}{lccl}
\toprule
\textbf{Disease} & $\boldsymbol{\rho_d}$ & $\boldsymbol{p_d}$ &
\textbf{Rationale for $p_d$} \\
\midrule
Dengue    & 0.3 & 0.5 & Mild/asymptomatic early phase; moderate travel
                        probability \\
Malaria   & 0.2 & 0.3 & Significant illness; partial immunity may allow
                        travel \\
Measles   & 0.6 & 0.05 & High symptomatic burden; travel while infectious
                        uncommon \\
Pertussis & 0.3 & 0.7 & Catarrhal stage resembles common cold; most
                        infected continue travel \\
\bottomrule
\end{tabular}
\end{table}

% ============================================================
\subsection*{NTTO 2026 visitor projections for top 12 source markets}

\begin{table}[h!]
\centering
\caption{NTTO 2026 international visitor projections and derived
  growth factors used in the WC-adjusted model. Source:
  US National Travel and Tourism Office Forecast Report
  (March 2025).}
\label{tab:ntto}
\begin{tabular}{lrrrc}
\toprule
\textbf{Country} & \textbf{2024 (thousands)} & \textbf{2026 (thousands)}
  & $\boldsymbol{\phi_c}$ & \textbf{WC team?} \\
\midrule
Canada         & 20,241 & 23,015 & 1.137 & Yes (host) \\
Mexico         & 16,990 & 19,557 & 1.151 & Yes (host) \\
United Kingdom & 4,037  & 4,467  & 1.107 & Yes (England, Scotland) \\
Japan          & 1,844  & 2,600  & 1.410 & Yes \\
China          & 1,626  & 2,641  & 1.624 & No \\
Germany        & 1,995  & 2,223  & 1.114 & Yes \\
Brazil         & 1,910  & 2,341  & 1.226 & Yes \\
South Korea    & 1,700  & 2,018  & 1.187 & Yes \\
India          & 2,190  & 2,541  & 1.161 & No \\
France         & 1,706  & 1,923  & 1.127 & Yes \\
Australia      & 1,025  & 1,244  & 1.214 & Yes \\
Italy          & 1,124  & 1,314  & 1.169 & No \\
\midrule
\textbf{Total international} & \textbf{72,390} & \textbf{85,017}
  & \textbf{1.174} & \\
\bottomrule
\end{tabular}
\end{table}

% ============================================================
\subsection*{Airport-level spatial allocation}

For the supplementary airport-level analysis, the expected
importations per city were further distributed across individual US
airports proportionally to annual enplanement volume. Enplanement
data were obtained from the Bureau of Transportation Statistics
Annual Airport Rankings (2023). Airports were spatially joined to US
counties using Census TIGER/Line cartographic boundary files (2023
vintage) via a point-in-polygon spatial join. Airport names were
matched to OurAirports global airport records using Jaro--Winkler
string distance with a state-constrained candidate set to handle
cross-state metropolitan areas (e.g., New York/New Jersey, Washington
DC/Virginia/Maryland). The probability that an international arrival
enters through airport $a$ was estimated as:
\begin{equation}
  \pi_a = \frac{E_a}{\sum_{a'} E_{a'}}
\end{equation}
where $E_a$ denotes annual enplanements at airport $a$ among the top
ranked US airports (excluding Alaska, Hawaii, and US territories).
Country-level expected importations were then allocated to airports
as $\lambda_{c,a,d} = \lambda_{c,d} \cdot \pi_a$.

% ============================================================
\subsection*{Sensitivity analysis}

All parameter values for $\rho_d$ and $p_d$ represent central
estimates subject to uncertainty. Given the lack of strong empirical
constraints on these values for large sporting events, a structured
sensitivity analysis varying $\rho_d$ and $p_d$ independently over a
plausible range (e.g., $\pm 50\%$ of central values) is recommended
before finalising results for policy communication.

% ============================================================
\subsection*{Software}

All analyses were performed in R (version 4.5.2). Key packages used
include \texttt{tidyverse} for data wrangling, \texttt{sf} and
\texttt{tigris} for spatial operations, \texttt{stringdist} for
fuzzy airport name matching, \texttt{readxl} for ingestion of Excel
data sources, and \texttt{cowplot} for figure composition. Code and
data are available at
\url{https://github.com/jdiestra1977/FIFA_worldCup_2026_risk}.

% ============================================================
\bibliographystyle{apalike}
\bibliography{references}

% ============================================================
\newpage
\appendix
\section*{Supplementary Table 1. All 48 nations qualified for the
2026 FIFA World Cup}
% ============================================================

\begin{longtable}{llc}
\toprule
\textbf{Country} & \textbf{Confederation} & \textbf{Host nation} \\
\midrule
\endfirsthead
\toprule
\textbf{Country} & \textbf{Confederation} & \textbf{Host nation} \\
\midrule
\endhead
\midrule
\multicolumn{3}{r}{\textit{Continued on next page}} \\
\endfoot
\bottomrule
\endlastfoot
% --- AFC (9) ---
Japan          & AFC & No \\
Iran           & AFC & No \\
South Korea    & AFC & No \\
Australia      & AFC & No \\
Qatar          & AFC & No \\
Saudi Arabia   & AFC & No \\
Uzbekistan     & AFC & No \\
Jordan         & AFC & No \\
Iraq           & AFC & No \\
% --- CAF (10) ---
Morocco        & CAF & No \\
Tunisia        & CAF & No \\
Egypt          & CAF & No \\
Algeria        & CAF & No \\
Ghana          & CAF & No \\
Cape Verde     & CAF & No \\
Senegal        & CAF & No \\
South Africa   & CAF & No \\
Ivory Coast    & CAF & No \\
DR Congo       & CAF & No \\
% --- CONCACAF (6) ---
Canada         & CONCACAF & Yes \\
Mexico         & CONCACAF & Yes \\
United States  & CONCACAF & Yes \\
Panama         & CONCACAF & No \\
Cura\c{c}ao   & CONCACAF & No \\
Haiti          & CONCACAF & No \\
% --- CONMEBOL (6) ---
Argentina      & CONMEBOL & No \\
Brazil         & CONMEBOL & No \\
Ecuador        & CONMEBOL & No \\
Paraguay       & CONMEBOL & No \\
Uruguay        & CONMEBOL & No \\
Colombia       & CONMEBOL & No \\
% --- UEFA (16) ---
England        & UEFA & No \\
France         & UEFA & No \\
Croatia        & UEFA & No \\
Portugal       & UEFA & No \\
Norway         & UEFA & No \\
Germany        & UEFA & No \\
Netherlands    & UEFA & No \\
Switzerland    & UEFA & No \\
Scotland       & UEFA & No \\
Spain          & UEFA & No \\
Austria        & UEFA & No \\
Belgium        & UEFA & No \\
Bosnia and Herzegovina & UEFA & No \\
Sweden         & UEFA & No \\
Turkey         & UEFA & No \\
Czech Republic & UEFA & No \\
% --- OFC (1) ---
New Zealand    & OFC & No \\
\end{longtable}

\noindent\small
\textit{Note}: Confederation abbreviations --- AFC: Asian Football
Confederation; CAF: Confederation of African Football; CONCACAF:
Confederation of North, Central America and Caribbean Association
Football; CONMEBOL: South American Football Confederation; UEFA:
Union of European Football Associations; OFC: Oceania Football
Confederation. Source: \url{https://en.wikipedia.org/wiki/2026_FIFA_World_Cup_qualification}
(accessed April 2026).

\end{document}
